\section{From a closed bridge to open sequents}
\label{sec:bridge}

A repository-wide framework should not force every reusable lemma to be closed
merely to move between the Foundation and Concise developments.  The audited
bridge is close to the desired result; the remaining work is an interface and
round-trip problem, not a new interpretation of arithmetic.

\subsection{The easy direction is already internally open}

The theorem
\code{toFoundation\_provable\_of\_axioms} has the form
\[
 \Prf\;Ax\;\Gamma\;p
 \;\longrightarrow\;
 \Gamma^{F,v}\;\vdash_{\Lambda}\;p^{F,v}
\]
for every assignment \(v\) of the \(n\) free Concise variables to Foundation
terms.  Its proof already handles hypotheses, all logical rules, equality
substitution, and binder lifting.  The public task is therefore:
\begin{enumerate}[label=(\alph*)]
\item specialize \(Ax\) to \(\PRAAx\) and discharge the already checked images
      of its axioms;
\item expose a theorem with a finite local context rather than immediately
      universally closing an empty-context result;
\item add rewrite lemmas showing compatibility with the repository's standard
      formula/context translation names.
\end{enumerate}
No new induction on \(\Prf\) is necessary.

\subsection{Reverse transport by finite deduction}

The existing \code{hilbertToConcise} translates an empty-context Hilbert proof
under an arbitrary assignment of free variables.  Let a Foundation
finite-context derivation have hypotheses
\(D_1,\ldots,D_r\) and conclusion \(P\).  Repeated use of the Foundation
finite-context deduction theorem gives an empty-context Hilbert proof of
\[
       D_1\to D_2\to\cdots\to D_r\to P.
\]
Translate this implication chain by \code{hilbertToConcise}.  In the Concise
calculus, reopen it with \code{Prf.impE} and the translated hypotheses.  This
yields the target
\[
  D_1^C,\ldots,D_r^C\quad\turnstile\quad P^C
\]
under every chosen free-term environment.

\begin{theorem}[Proposed open Foundation-to-Concise transport]
\label{thm:open-bridge}
For the reference Foundation iPRA theory and every finite context \(\Delta\),
free-variable environment \(v\), and formula \(P\), a Foundation derivation
\(\Delta\vdash P\) induces
\[
      \Delta^{C,v}\turnstile P^{C,v}.
\]
\end{theorem}

\begin{proof}[Detailed implementation proof]
Induct on the finite list \(\Delta\), using Foundation's \code{deduct} theorem
to move its head into an implication until the context is empty.  Apply the
existing Hilbert-to-Concise translator with the same environment \(v\).  The
translation of implication is definitionally the Concise implication.  Now
induct over \(\Delta\) in the opposite order and apply modus ponens with
\code{Prf.hyp} at the corresponding list membership.  If the Foundation
finite-context representation stores a proof object rather than only a
provability proposition, the same construction is a direct recursion on that
object.  Binder cases are already handled by \code{underBinder} and the
proved image lemmas in \code{FoundationToConcise.lean}.
\end{proof}

\subsection{The free-variable environment}

A closed theorem can assign every free Foundation variable to zero because
there are none in the sentence.  An open reflection theorem cannot.  For a
formula whose active free variables are \(0,\ldots,n-1\), choose the Concise
environment
\[
        v(i)=\begin{cases}
          \PTm.\mathsf{var}(i),&i<n,\\
          0,&i\ge n.
        \end{cases}
\]
and prove that the translation depends only on the active finite support.
The analogous Foundation environment uses its first \(n\) free variables.
Substitution compatibility then gives transport at arbitrary terms by
specializing the open theorem, rather than re-running the bridge.

\subsection{Round trips and reflection}

Preservation in both directions does not by itself give an open-sequent
\emph{equivalence}.  One also needs formula round trips.

\begin{description}
\item[Concise round trip.]  The current bridge reports that the
Concise-to-Foundation-to-Concise translation is exact.  Expose this equality
for open formulas and show it commutes with context maps.
\item[Foundation round trip.]  The reverse round trip can replace a
positive-arity zero symbol by a composition of the nullary zero.  Prove these
terms equal in the Foundation theory, then prove formulas equivalent by
structural induction and equality substitution.
\end{description}

With local equivalences \(D\leftrightarrow D^{FC}\) for every hypothesis and
\(P\leftrightarrow P^{FC}\) for the goal, use simultaneous cut to replace the
round-tripped context.  This yields the desired open-sequent reflection theorem
up to the already documented zero-code normalization.

Primitive-recursive function codes require a structural induction.  In the
primitive-recursion constructor, compare the two interpreted iterators by qf
induction on the recursion argument.  This is exactly the kind of theorem the
bridge's closed API previously avoided but the ordinary iPRA core now supports.

\subsection{Bridge acceptance tests}

A convincing integration should include more than a closed sentence.
\begin{enumerate}[label=(B\arabic*)]
\item a two-variable equality theorem under a local equality hypothesis;
\item substitution of non-variable arithmetic terms after transport;
\item a theorem with a quantifier and one shifted local hypothesis;
\item a qf-induction theorem with parameters;
\item both round trips on a primitive-recursive symbol containing composition
      and recursion;
\item an SCR contract transported with its certificate hypothesis still open.
\end{enumerate}
The last test is important: the open bridge should preserve the very local
producer boundaries that make the framework modular.
